% Slide template for the FreeBSD Developer Summit.
% Take it or leave it :-)
%
% It requires LaTeX and LaTeX Beamer [1] to compile.
% pdfLaTeX is recommended for compilation as it produces
% PDF file immediately.
%
% $ pdflatex some.latex
%
%
% It is also recommended to convert images to PDF
% by using ImageMagick ("convert") before including them.
%
% [1] http://latex-beamer.sourceforge.net
%

\documentclass{beamer}

\usepackage{url}
%\usepackage[T1,T2A]{fontenc}
\usepackage[T1]{fontenc}
\usepackage[utf8]{inputenc}
\usepackage{newunicodechar}
\usepackage{textcomp}
%\usepackage[english,russian]{babel}
\usepackage[english]{babel}
\usepackage{verbatim}
\usepackage{graphicx}
\usepackage{listings}
\usepackage{mathtools}
\usepackage{color}
\usepackage{listings}
\usepackage{tikz}

\mode<presentation>
{
  \definecolor{beamer@gker}{rgb}{0.0,0.47,0.44}
  \setbeamercolor*{structure}{fg=beamer@gker}
  \logo{\includegraphics[scale=0.13]{../ganditemplate/Gandi_logo_black.eps}}
}

\setbeamertemplate{footline}[text line]{%
  \parbox{\linewidth}{\vspace*{-8pt}
  \prestitle
  \hfill
  \insertshorttitle
  \hfill
  \insertframenumber\ of \inserttotalframenumber}}
\setbeamertemplate{navigation symbols}{}

\definecolor{light-gray}{gray}{0.60}

\lstdefinestyle{custommake}{
  belowcaptionskip=1\baselineskip,
  breaklines=true,
  frame=L,
  xleftmargin=\parindent,
  language=Make,
  showstringspaces=false,
  basicstyle=\footnotesize\ttfamily,
  keywordstyle=\bfseries\color{green!40!black},
  commentstyle=\itshape\color{purple!40!black},
  identifierstyle=\color{blue},
  stringstyle=\color{orange},
}

\definecolor{mygreen}{rgb}{0,0.6,0}
\definecolor{mygray}{rgb}{0.5,0.5,0.5}
\definecolor{mymauve}{rgb}{0.58,0,0.82}



\lstset{ %
  backgroundcolor=\color{white},   % choose the background color; you must add \usepackage{color} or \usepackage{xcolor}
  basicstyle=\ttfamily\tiny,        % the size of the fonts that are used for the code
  breakatwhitespace=false,         % sets if automatic breaks should only happen at whitespace
  breaklines=true,                 % sets automatic line breaking
  captionpos=b,                    % sets the caption-position to bottom
  commentstyle=\color{mygreen},    % comment style
  deletekeywords={...},            % if you want to delete keywords from the given language
  escapeinside={\%*}{*)},          % if you want to add LaTeX within your code
  extendedchars=true,              % lets you use non-ASCII characters; for 8-bits encodings only, does not work with UTF-8
  frame=single,                    % adds a frame around the code
  keepspaces=true,                 % keeps spaces in text, useful for keeping indentation of code (possibly needs columns=flexible)
  keywordstyle=\color{blue},       % keyword style
  language=Make,                 % the language of the code
  morekeywords={*,...,.include},            % if you want to add more keywords to the set
  numbers=left,                    % where to put the line-numbers; possible values are (none, left, right)
  numbersep=5pt,                   % how far the line-numbers are from the code
  numberstyle=\tiny\color{mygray}, % the style that is used for the line-numbers
  rulecolor=\color{black},         % if not set, the frame-color may be changed on line-breaks within not-black text (e.g. comments (green here))
  showspaces=false,                % show spaces everywhere adding particular underscores; it overrides 'showstringspaces'
  showstringspaces=false,          % underline spaces within strings only
  showtabs=false,                  % show tabs within strings adding particular underscores
  stepnumber=1,                    % the step between two line-numbers. If it's 1, each line will be numbered
  stringstyle=\color{mymauve},     % string literal style
  tabsize=8,                       % sets default tabsize to 2 spaces
  title=\lstname                   % show the filename of files included with \lstinputlisting; also try caption instead of title
}

%\addtobeamertemplate{frametitle}{}{%
%\begin{tikzpicture}[remember picture,overlay]
%	\node[anchor=north east,yshift=-5pt] at (current page.north east) {\includegraphics[scale=0.13]{../ganditemplate/Gandi_logo_black.eps}};
%\end{tikzpicture}}

%%\lstset{escapechar=@,style=custommake}

\usepackage{color}
\usepackage{listings}
\usepackage{tipa}
\usepackage{svg}
\definecolor{indigo}{rgb}{0.29, 0.0, 0.51}
\newcommand{\prestitle}{sysadmin\#6}
\newcommand{\prev}[1]{\href{https://svnweb.freebsd.org/changeset/ports/#1}{r#1}}
\lstset{
	numbers=none,
	frame=l,
	aboveskip=0pt, belowskip=0pt
}

\usepackage{tikz}

\usepackage{wordcloud}

\usepackage{xcolor} % for \textcolor
\usepackage{xparse} % for \DeclareExpandableDocumentCommand
\ExplSyntaxOn
\DeclareExpandableDocumentCommand { \fpeval } { m } { \fp_eval:n {#1} }
\ExplSyntaxOff

\newcommand{\cloudentry}[2]{
\renewcommand{\magstep}{#2}
\larger #1\smaller}

\title{Gestion de services sous Unix}
\author{Baptiste Daroussin \\ \url{bapt@FreeBSD.org} \\ \url{bapt@gandi.net}}
\institute{\includegraphics[scale=0.13]{../gandi/Gandi_logo_black.eps}}
\date{sysadmin \#6 \\ Paris \\ 19 Février 2016}

\begin{document}
\begin{frame}[plain]
	\titlepage
\end{frame}

\begin{frame}
	\frametitle{Disclaimer}
	\center
	\includegraphics[scale=0.20]{trollfreezone.eps}
\end{frame}

\begin{frame}[fragile]
	\frametitle{L'arlésienne}
	\center
	\wordcloud[number-of-words = 33,
		style = {\fboxsep=1pt\fboxrule=0pt%
		\textcolor[rgb]{\fpeval{1-#2/\wordcloudMaxValue} .5 .5}{#1}} ,
		angle-step = 15
	]{
		rcNG launchd SysVinit SystemXVI SMF relaunchd nosh systemd OpenRC
		daemontools upstart monit Epoch finit Runit initng rc minit sinit
		ninit perp watchman sparkrc launchd SMF systemd systemd OpenRC rcNG rcNG
		supervisor svc circus Shepherd procd DEMONS eINIT
		OpenRC svc daemontools SMF watchman relaunchd rc eINIT systemd
	}
\end{frame}

\begin{frame}
	\frametitle{Il était une fois init(8)}
	\begin{itemize}
		\item petit mais costaud
			\pause
		\item I'm your father(c)(tm)
			\pause
		\item laisse la main a d'autres pour la gestion des services... \pause presque
			\pause
			\begin{itemize}
				\item /etc/inittab
				\item /etc/ttys
			\end{itemize}
			\pause
		\item Généralement des scripts shell
			\begin{itemize}
				\item SysVinit
				\item rcNG
				\item rc
			\end{itemize}
	\end{itemize}
\end{frame}

\begin{frame}
	\frametitle{Alors pourquoi tout ce bordel?}
	\begin{itemize}
			\pause
		\item trop naïf
			\pause
		\item framework complexes
			\pause
		\item cohérence compliquée
			\pause
		\item gestion de ressource complexe
			\pause
		\item peu réactif
	\end{itemize}
\end{frame}

\begin{frame}
	\frametitle{Trop naïf}
	\begin{itemize}
			\pause
		\item ne peut superviser qu'un process simple (via son pid)
			\pause
		\item perdu si double fork
			\pause
		\item interactions limitées
			\pause
		\item gestion de dépendance faible
	\end{itemize}
\end{frame}

\begin{frame}
	\frametitle{Framework complexes}
	\begin{itemize}
			\pause
		\item généralement en shell
			\pause
		\item complexe à isoler (cf commande service(8))
			\pause
		\item complexe à garder une cohérence entre les scripts d'init
			\pause
		\item complexe d'intégrer les systèmes de gestion de ressources
			\pause
		\item souvent inefficace (performances)
			\pause
		\item pas de mécanismes de surveillance de processus
			\pause
		\item peu de contrôle sur les processus (double fork)
			\pause
		\item fiabilité aléatoire du statut
	\end{itemize}
\end{frame}

\begin{frame}
	\frametitle{Réactivité}
	\begin{itemize}
			\pause
		\item gestion événementielle forcément externe (cron, udev/devd, inetd)
			\pause
			\begin{itemize}
				\item duplication de la gestion des services
				\item incohérence dans la gestion des services
			\end{itemize}
			\pause
		\item gestion événementielle souvent peu réactive: latence entre l'arrivée de l'évènement et l'action
			\pause
		\item fault management absent
	\end{itemize}
\end{frame}

\begin{frame}
	\frametitle{D'autre soucis: pas simple d'être un bon daemoniste}
	\begin{itemize}
			\pause
		\item utiliser setuid(2), setgid(2) ET setgroup(2) (Vérifier que l'opération c'est réellement effectuée)
			\pause
		\item toujours faire un chdir(2) après un chroot(2) (Vérifier leur bon fonctionnement)
			\pause
		\item configurer de manière portable les setrlimit(2)
			\pause
		\item utiliser les mécanismes OS spécifiques de sandboxing
	\end{itemize}
\end{frame}

\begin{frame}
	\frametitle{Le problème des PIDs}
	\begin{itemize}
			\pause
		\item concept "racy"
			\pause
		\item services avec multiple workers hors de contrôle :
			\begin{itemize}
				\item les fils qui échappent à la surveillance des parents
				\item les parents laxistes qui abandonnent leurs enfants
			\end{itemize}
	\end{itemize}
\end{frame}

\begin{frame}
	\frametitle{Ça manque de sucre un peu tout ça}
	\begin{itemize}
		\item y a des trous dans ma syslog :(
		\item c'est pas très parallèle tout ça
		\item comment je m'assure que ça marche?
	\end{itemize}
\end{frame}

\begin{frame}
	\frametitle{Des solutions possibles?}
	\begin{itemize}
		\item Pour les PIDs :
			\begin{itemize}
					\pause
				\item process descriptors (pas suffisant) (Linux 4.1, FreeBSD 10.0)
					\pause
				\item les reapers (Linux 3.4, FreeBSD 10.2, Dragonfly 4.0, Solaris 10)
			\end{itemize}
			\pause
		\item Gérer les ressources le sandboxing et la daemonisation via de génération de services
			\pause
		\item La réactivité : faire communiquer nativement les différents composants
			\pause
		\item La complexité : décrire au maximum les services via des format de description
	\end{itemize}
\end{frame}

\begin{frame}
	\frametitle{Et un peu de sucre en poudre}
	\begin{itemize}
			\pause
		\item Préparer les sockets :
			\begin{itemize}
				\item aide à la parallélisation
				\item simplifie la séparation de privilèges
				\item aide à la robustesse
			\end{itemize}
			\pause
		\item utiliser un IPC:
			\begin{itemize}
				\item limite la duplication de code
				\item permet la séparation des tâches (KISS)
			\end{itemize}
	\end{itemize}
\end{frame}

\begin{frame}
	\frametitle{Oui mais...}
	\begin{itemize}
			\pause
		\item Faut il remplacer init(8) ? \pause pas nécessairement
			\pause
		\item Un des systèmes existant permet il corrige t il tout ça ? \pause NON
			\pause
		\item Un des systèmes adresse t il les problèmes majeures ? \pause OUI
			\pause
		\item Allez lequel ?\pause
	\end{itemize}
	\center
	\huge SMF \\
	\pause
	\normalsize
	mais over engineer :(
\end{frame}

\begin{frame}
	\frametitle{Un peu de lecture}
	\begin{itemize}
		\item \href{https://www.usenix.org/legacy/event/lisa05/tech/full_papers/adams/adams.pdf}{Solaris Service Management Facility: Modern System Startup and Administration}
		\item \href{https://www.usenix.org/legacy/event/usenix01/freenix01/full_papers/mewburn/mewburn_html/}{The Design and Implementation of the NetBSD rc.d system}
		\item \href{https://svnweb.freebsd.org/base/head/usr.sbin/cron/cron/do_command.c?view=markup}{Les commentaires de Paul Vixie dans do\_command.c}
	\end{itemize}
\end{frame}

\begin{frame}
	\frametitle{Questions?}
	\center
	\includegraphics[scale=0.25]{dont-feed-trolls.eps}
\end{frame}

\end{document}
